\documentclass[aspectratio=169,xcolor=dvipsnames]{beamer}

\usepackage{amsthm}
\usepackage{amsmath}
\usepackage{amssymb}
\usepackage{tikz}
\usepackage{pgfplots}
\usepackage{xfp}
\pgfplotsset{compat=1.18}
\usepackage[authoryear, round]{natbib}
\usepackage{hyperref}
\usetikzlibrary{arrows.meta,positioning,decorations.pathreplacing}
\hypersetup{
    colorlinks=true,
    linkcolor=blue,
    citecolor=blue,
    urlcolor=blue
}

\definecolor{darkgreen}{rgb}{0,0.5,0}
\definecolor{darkblue}{rgb}{0,0,0.55}

\newtheorem{proposition}{Proposition}
\setbeamertemplate{footline}[frame number]
\usetheme{Frankfurt}
\usefonttheme{serif}

\title{ECON8037: Financial Economics}
\subtitle{\normalsize Week 6 Lecture - Liquidity and Asset Prices}
\author{Simon Mishricky}
\institute{Research School of Economics, Australian National University}
\date{\today}

\begin{document}

\section{Lagos with Two Assets}
\frame{\titlepage}

\begin{frame}{Lagos}
    Continuing with \textcolor{blue}{Lagos (2010)}, we extend the model to situations in which bonds and shares have different (exogenous) liquidity properties. Specifically, suppose an agent can find himself in two types of meetings in the decentralised market: With probability $\theta_2$, he is in a meeting where he can use any of the two assets for payment, which with probability $\theta_1 = 1 - \theta_2$, he can only use bonds. Thus $\theta_1 \in [0, 1]$ indexes the degree of ``illiquidity'' of equity shares. (The subscript refers to the number of assets that can be used for payment in that particular type of meeting).
\end{frame}

\begin{frame}{Lagos}
    The value function of an agent who enters the centralised market with portfolio $\mathbf{a}_t$ in a period when dividends are $d_t$ still satisfies
    \[
        W(\mathbf{a}_t, d_t) = \max_{c_t, n_t, h_t, \mathbf{a}_t} \{U(c_t) + v(n_t) - A h_t + \beta \mathbb{E} V(\mathbf{a}_{t+1}, d_{t+1})\}
    \]
    Subject to
    \[
        c_t + w_t n_t + \boldsymbol{\phi}_t \mathbf{a}_{t+1} = (\phi^s_t + d_t) a^s_t + a^b_t + w_t h_t - \tau_t
    \]
    \[
        0 \le c_t, \quad 0 \le n_t, \quad 0 \le h_t \le \bar n, \quad 0 \le \mathbf{a}_{t+1}
    \]
\end{frame}

\begin{frame}{Lagos}
    As before, the terms of trade in each meeting will be determined by bargaining. Since there are now two types of meetings, there will be two bargaining solutions. For simplicity, we will refer to a meeting where $i$ assets can be used for payment as a ``meeting of type $i$''.

    The terms of trade in a meeting of type $i$ between a buyer who holds portfolio $\mathbf{a}_t = (a^s_t, a^b_t)$ and a seller with portfolio $\tilde{\mathbf{a}}_t = (\tilde a^s_t, \tilde a^b_t)$ are $[q^i(\mathbf{a}_t, \tilde{\mathbf{a}}_t), \mathbf{p}^i(\mathbf{a}_t, \tilde{\mathbf{a}}_t)]$, where $q^i(\mathbf{a}_t, \tilde{\mathbf{a}}_t) \in \mathbb{R}_+$ and $\mathbf{p}^i(\mathbf{a}_t, \tilde{\mathbf{a}}_t) \in \mathbb{R}^2_+$.

    In particular, $\mathbf{p}^2(\mathbf{a}_t, \tilde{\mathbf{a}}_t) = [p^s(\mathbf{a}_t, \tilde{\mathbf{a}}_t), p^b(\mathbf{a}_t, \tilde{\mathbf{a}}_t)]$, where $p^s(\mathbf{a}_t, \tilde{\mathbf{a}}_t)$ denotes the quantity of shares, and $p^b(\mathbf{a}_t, \tilde{\mathbf{a}}_t)$ the quantity of bonds that the buyer hands over to the seller. The trading restrictions imply that $\mathbf{p}^1(\mathbf{a}, \tilde{\mathbf{a}}_t) = [0, p^b(\mathbf{a}_t, \tilde{\mathbf{a}}_t)]$.
\end{frame}

\begin{frame}{Lagos}
    The value of an agent who enters the decentralised market holding portfolio $\mathbf{a}$ in a period when the dividend realisation is $s$, satisfies
    \begin{align*}
        V(\mathbf{a}, s) &= \alpha \sum_{i=1}^2 \theta_i \int \{u[q^i(\mathbf{a}, \tilde{\mathbf{a}})] + W[\mathbf{a} - \mathbf{p}^i(\mathbf{a}, \tilde{\mathbf{a}}), s]\} d\mathbf{G}(\tilde{\mathbf{a}}) \\
        &\quad + \alpha \sum_{i=1}^2 \theta_i \int \{-e[q^i(\tilde{\mathbf{a}}, \mathbf{a})] + W[\mathbf{a} + \mathbf{p}^i(\tilde{\mathbf{a}}, \mathbf{a}), s]\} d\mathbf{G}(\tilde{\mathbf{a}}) + (1 - 2\alpha) W(\mathbf{a}, s)
    \end{align*}
    Where $\mathbf{G}$ denotes the distribution of portfolios.
\end{frame}

\begin{frame}{Lagos}
    The terms of trade $[q^i(\mathbf{a}, \tilde{\mathbf{a}}), \mathbf{p}^i(\mathbf{a}, \tilde{\mathbf{a}})]$ for $i = 1, 2$, are still determined by Nash bargaining where the buyer has all the bargaining power. Now let $\boldsymbol{\lambda}^1_t = (0, \lambda^b_t)$, where $\lambda^b_t = \frac{A}{w_t}$, and $\boldsymbol{\lambda}^2_t = \boldsymbol{\lambda}_t = (\lambda^s_t, \lambda^b_t)$, where $\lambda^s_t = \frac{A}{w_t}(\phi^s_t + d_t)$. Then, the bargaining solution is $q^i(\mathbf{a}, \tilde{\mathbf{a}}_t) = q(\boldsymbol{\lambda}^i_t \mathbf{a}_t)$, where the function $q(\cdot)$ is still given by
    \[
        q(\boldsymbol{\lambda}_t \mathbf{a}_t) = \begin{cases} q^* & \text{if } \boldsymbol{\lambda}_t \mathbf{a}_t \ge e(q^*) \\ e^{-1}(\boldsymbol{\lambda}_t \mathbf{a}_t) & \text{if } \boldsymbol{\lambda}_t \mathbf{a}_t < e(q^*) \end{cases}
    \]
\end{frame}

\begin{frame}{Lagos}
    Using the bargaining solution and the fact that $W[\mathbf{a} + \mathbf{p}^i(\mathbf{a}, \tilde{\mathbf{a}}), d] - W(\mathbf{a}, d) = \boldsymbol{\lambda}_t \mathbf{p}^i(\mathbf{a}, \tilde{\mathbf{a}})$, the value of search can be written as:
    \[
        V(\mathbf{a}, d) = \sum_{i=1}^2 \theta_i \mathcal{S}(\boldsymbol{\lambda}^i \mathbf{a}) + W(\mathbf{a}, d)
    \]
    In the centralised market, the agent's choices of $c_t$ and $n_t$ are still characterised by
    \[
        U'(c_t) = \frac{A}{w_t}
    \]
    \[
        v'(n_t) = A
    \]
\end{frame}

\begin{frame}{Lagos}
    The first-order conditions for asset holdings $\mathbf{a}_{t+1}$ are
    \[
        \frac{A}{w_t} \phi^i_t = \beta \mathbb{E} \frac{dV(\mathbf{a}_{t+1}, d_{t+1})}{d a^i_{t+1}}, \quad i = s, b
    \]
    Where
    \begin{align*}
        \frac{dV(\mathbf{a}_{t+1}, d_{t+1})}{d a^b_{t+1}} &= \left[ 1 + \alpha \sum_{i=1}^2 \theta_i \left( \frac{u'[q(\boldsymbol{\lambda}^i_{t+1} \mathbf{a}_{t+1})]}{e'[q(\boldsymbol{\lambda}^i_{t+1} \mathbf{a}_{t+1})]} - 1\right)\right] \lambda^b_{t+1} \\
        \frac{dV(\mathbf{a}_{t+1}, d_{t+1})}{d a^s_{t+1}} &= \left[ 1 + \alpha \theta_2 \left( \frac{u'[q(\boldsymbol{\lambda}^2_{t+1} \mathbf{a}_{t+1})]}{e'[q(\boldsymbol{\lambda}^2_{t+1} \mathbf{a}_{t+1})]} - 1\right)\right] \lambda^s_{t+1}
    \end{align*}
    As usual, the agent's choices do not depend on his individual asset holdings, and the distribution of assets will be degenerate in equilibrium.
\end{frame}

\begin{frame}{Lagos}
    Given the dividend process $\{d_t\}_{t=0}^\infty$ and a path $\{B_t, \tau_t\}_{t=0}^\infty$, an equilibrium is an allocation $\{c_t, n_t, \mathbf{a}_{t+1}\}_{t=0}^\infty$ together with the set prices $\{w_t, \boldsymbol{\phi}_t\}_{t=0}^\infty$ and bilateral terms of trade $\{q_t\}_{t=0}^\infty$, such that:
    \begin{enumerate}[(a)]
        \item The individual choices $\{c_t, n_t, \mathbf{a}_{t+1}\}_{t=0}^\infty$ solve the agent's problem in the decentralised market, given prices;
        \item The terms of trade are determined by Nash bargaining, i.e.,
        \[
            q^1_t = \min \left\{ q^*, e^{-1}\left( \frac{A}{w_t} a^b_t\right)\right\} \quad \text{and} \quad q^2_t = \min \left\{ q^*, e^{-1}\left[ \frac{A}{w_t}(\phi^s_t + d_t) a^s_t + \frac{A}{w_t} a^b_t\right]\right\}
        \]
        \item Prices are such that the centralised market clears, i.e., $c_t = d_t$, $a^s_{t+1} = 1$, and $a^b_{t+1} = B_{t+1}$.
    \end{enumerate}
\end{frame}

\section{Recursive Equilibrium}

\begin{frame}{Lagos}
    In a \emph{recursive equilibrium}, allocations and prices are time-invariant of the aggregate state $x$, the current dividend realisation. Specifically, $\phi^s_t = \phi^s(x_t)$, $\phi^b_t = \phi^b(x_t)$ and $w_t = w(x_t) \equiv \frac{A}{U'(x_t)}$. In addition, $\lambda^s_t = \lambda^s(x_t) \equiv U'(x_t)[\phi^s(x_t) + x_t]$, and $\lambda^b_t = U'(x_t)$. Also, suppose that the government follows the stationary policy $B_{t+1} = B(x_t)$, and that the aggregate state follows a Markov process with $\mathbb{P}\{x_{t+1} \le x' | x_t = x\} = F(x', x)$.
\end{frame}

\begin{frame}{Lagos}
    Then, asset prices satisfy:
    \[
        \lambda^s(x) - U'(x) x = \beta \int L^s(x') \lambda^s(x') dF(x', x)
    \]
    \[
        U'(x) \phi^b(x) = \beta \int L^b(x') U'(x') dF(x', x)
    \]
    Where
    \[
        L^s(x') = 1 + \alpha(1 - \theta) \left( \frac{u'\{q[\lambda^s(x') + U'(x') B']\}}{e'\{q[\lambda^s(x') + U'(x') B']\}} - 1\right)
    \]
    $B' = B(x)$, and
    \[
        L^b(x') = L^s(x') + \alpha \theta \left( \frac{u'\{q[U'(x') B']\}}{e'\{q[U'(x') B']\}} - 1\right)
    \]
    To simplify notation, we have set $\theta_1 = \theta$.
\end{frame}

\section{Equity Premium and Risk-Free Rate Puzzles}

\begin{frame}{Lagos}
    \textbf{The Equity Premium and Risk-Free Rate Puzzles}

    Use the \emph{intrinsic returns} $\hat R^b(x)$ and $\hat R^s(x, x')$, to define the \emph{full returns} $R^b(x, x') = L^b(x') \hat R^b(x)$ and $R^s(x, x') = L^s(x') \hat R^s(x, x')$ and write the asset prices as
    \[
        \int \beta \frac{U'(x')}{U'(x)} R^i(x, x') dF(x', x) = 1
    \]
    for $i = s, b$ and all $x$.
\end{frame}

\begin{frame}{Lagos}
    Let $\bar F$ be the invariant distribution for the transition function $F$; then the Euler equations lead to:
    \[
        \int \int \beta \frac{U'(x')}{U'(x)}[R^s(x, x') - R^b(x, x')] dF(x', x) d\bar F(x) = 0
    \]
    \[
        \int \int \beta \frac{U'(x')}{U'(x)} R^b(x, x') dF(x', x) d\bar F(x) - 1 = 0
    \]
    a pair of statistical restrictions on equilibrium asset returns $R^i(x, x')$, and the marginal rate of substitution, $\beta \frac{U'(x')}{U'(x)}$.
\end{frame}

\begin{frame}{Lagos}
    For the special case of an economy with no liquidity needs (i.e., $\alpha = 0$), the intrinsic returns equal the full returns, i.e., $\hat R^s(x, x') = R^s(x, x')$ and $\hat R^b(x) = R^b(x, x')$. In this case, one can estimate the expectations of the left hand sides of the Euler equations using the sample means
    \[
        \hat \omega^e = \frac{1}{T} \sum_{t=1}^T \beta \frac{U'(c_{t+1})}{U'(c_t)} [\hat R^s_{t+1} - \hat R^b_t]
    \]
    \[
        \hat \omega^b = \frac{1}{T} \sum_{t=1}^T \beta \frac{U'(c_{t+1})}{U'(c_t)} \hat R^b_t - 1
    \]
    A vast body of work has been devoted to trying to rationalise the finding that for standard parameterisation of preferences, the statistical restrictions $\hat \omega^e = \hat \omega^b = 0$ are violated by U.S. data. The finding that $\hat \omega^e > 0$ constitutes the \textcolor{blue}{equity premium puzzle}, while $\hat \omega^b < 0$ is referred to as the \textcolor{blue}{risk-free rate puzzle}.
\end{frame}

\begin{frame}{Lagos}
    The statistics $\hat \omega^e$ and $\hat \omega^b$ that define these puzzles are constructed using only the \emph{intrinsic returns}, $\hat R^s_{t+1}$ and $\hat R^b_t$. But according to the model that we have developed, agents price assets using the \emph{full returns} $R^s_{t+1}$ and $R^b_{t+1}$, that is, the intrinsic returns augmented by their respective liquidity factors.
\end{frame}

\begin{frame}{Lagos}
    In our case we can write
    \[
        \int \int \beta \frac{U'(x')}{U'(x)} [\hat R^s(x, x') - \hat R^b(x, x')] dF(x', x) d\bar F(x) = \omega^e
    \]
    \[
        \int \int \beta \frac{U'(x')}{U'(x)} \hat R^b(x, x') dF(x', x) d\bar F(x) - 1 = \omega^b
    \]
    Where
    \[
        \omega^e = \int \int \beta \frac{U'(x')}{U'(x)} \left\{ [L^b(x') - 1] \hat R^b(x) - [L^s(x') - 1] \hat R^s(x, x')\right\} dF(x, x') d\bar F(x)
    \]
    \[
        \omega^b = -\int \int \beta \frac{U'(x')}{U(x)} [L^b(x') - 1] \hat R^b(x) dF(x, x') d\bar F(x)
    \]
\end{frame}

\begin{frame}{Lagos}
    Note that $\omega^e$ and $\omega^b$ are the theoretical counterparts to $\hat \omega^e$ and $\hat \omega^b$, respectively. The left hand side of the first Euler equation is a weighted average of the excess intrinsic return of equity over bonds, which in equilibrium must equal $\omega^e$. The wedge $\omega^e$ will be larger, the larger $\frac{L^b - 1}{L^s - 1} - \frac{\hat R^s}{\hat R^b}$ is on average, i.e., the larger is the excess liquidity return of bonds over equity relative to the excess intrinsic return of equity over bonds.
\end{frame}

\begin{frame}{Lagos}
    The sign of $\omega^e$ is ambiguous in general. But to the extent that bonds are more liquid than equity (in the sense of larger $\theta$), the latter will tend to pay an ``illiquidity premium'', implying $\omega^e > 0$. Note that even without assuming any exogenous liquidity differences between equity and bonds, we have $\omega^b < 0$, to the extent that there are liquidity needs at least for some realisations of the dividend process. Thus, qualitatively, the model with liquidity needs is always consistent with the risk-free rate puzzle. And in addition, if the liquidity return on bonds is large enough relative to that on equity, then the model may also help rationalise the equity premium puzzle.
\end{frame}

\begin{frame}[allowframebreaks]{References}
    \small
    \begin{thebibliography}{99}
        \bibitem[Lagos(2010)]{lagos2010}
            Lagos, R.\ Asset prices and liquidity in an exchange economy. \textit{Journal of Monetary Economics}, 57(8):913--930, 2010.
    \end{thebibliography}
\end{frame}

\end{document}
